\chapter{Modelo del mecanismo}
\label{chap:modelo}

Este capítulo desarrolla el modelo estático del mecanismo: relaciona el giro del tambor con el
ángulo de la barra, la longitud del cable y su tensión, y de ahí obtiene el par. Con ese modelo se
comparan cuatro perfiles de tambor y se establece el método con el que, en el
capítulo~\ref{chap:resultados}, se contrastan las medidas. Las mismas ecuaciones están
implementadas en \fichero{static/mecanismo.js}, de modo que el simulador y el visor de ensayos
calculan exactamente lo mismo que aquí se escribe.

El resultado central del capítulo es la sección~\ref{sec:mod-basta}: un cilindro con el eje
desplazado —un tambor \emph{excéntrico}— reproduce el perfil de par constante con un error de
menos del \SI{1}{\percent}, sin necesidad de mecanizar una espiral.

\section{Descripción del mecanismo}

La figura~\ref{fig:mecanismo} muestra el esquema. Una barra rígida articulada en O lleva suspendida
una masa $m$ a una distancia $L_3+L_4$ de la articulación. El cable se ancla a la barra en el punto
A, a una distancia $L_3$ de O, y sube hasta el punto de salida D, desde donde se enrolla en el
tambor. El tambor gira solidario al eje de salida de la transmisión, con una reducción $N$ de
vueltas de motor por vuelta de tambor.

El punto D \emph{no} está sobre la vertical de O: está a una altura $L_2$ y desplazado $dx$ hacia el
lado contrario a la barra. Ese detalle, que en una primera versión del modelo se pasó por alto,
cambia la física y no sólo el dibujo, como se ve en la sección~\ref{sec:mod-tension}.

Como variable angular se usa $\alpha$, el ángulo en O entre la línea O$\to$D y la barra, que recorre
de 0° (barra plegada sobre esa línea) a 180° (barra en su prolongación). El ángulo respecto a la
horizontal es $\theta = \beta - \alpha$, siendo $\beta$ la inclinación de O$\to$D. Se prefiere
$\alpha$ porque a cada longitud de cable le corresponde un único $\alpha$, sin la ambigüedad que
tenía medir desde la horizontal.

\begin{figure}[H]
\centering
\begin{tikzpicture}[scale=1.7,>=Stealth]
  \def\al{100}
  \coordinate (O) at (0,0);
  \coordinate (D) at (-0.65,2.9);
  \coordinate (A) at ({0.95*cos(12)},{0.95*sin(12)});
  \coordinate (W) at ({2.05*cos(12)},{2.05*sin(12)});
  \fill[gray!25] (-0.35,-0.12) rectangle (0.35,0);
  \draw[gray!55,line width=1.2pt,dashed] (O) -- (D);
  \draw[dashed,gray] (O) -- (1.3,0);
  \draw[line width=2.4pt] (O) -- (W);
  \draw[blue!70!black,thick] (D) -- (A) node[pos=0.45,right] {$c$};
  \draw[violet,thick,fill=violet!8] (-0.348,3.076) circle (0.35);
  \fill[violet] (-0.348,3.076) circle (0.9pt);
  \node[violet,right,font=\small] at (0.08,3.28) {tambor $r(\varphi)$};
  \node[below right=-1pt,font=\small] at (-0.52,2.70) {D};
  \fill (D) circle (1.2pt);
  \fill (O) circle (1.6pt) node[below=3pt] {O};
  \fill[blue!70!black] (A) circle (1.4pt) node[below right=-1pt] {A};
  \draw[->] (0.45,0) arc (0:12:0.45);
  \node[font=\small] at (0.62,0.07) {$\theta$};
  \draw[->,green!45!black] ({0.33*cos(\al+12)},{0.33*sin(\al+12)}) arc (\al+12:12:0.33);
  \node[green!45!black,font=\small] at (0.02,0.46) {$\alpha$};
  \draw (W) -- ++(0,-0.38);
  \fill[orange!80!black] ($(W)+(-0.15,-0.66)$) rectangle ($(W)+(0.15,-0.38)$);
  \draw[->,thick] ($(W)+(0,-0.7)$) -- ++(0,-0.4) node[right] {$mg$};
  \draw[->,thick,red!75!black] (A) -- ($(A)!0.3!(D)$) node[left=2pt] {$T$};
  \draw[<->] (-1.15,0) -- (-1.15,2.9) node[midway,left] {$L_2$};
  \draw[dotted] (-1.2,2.9) -- (-0.7,2.9);
  \draw[<->] (-0.65,-0.35) -- (0,-0.35) node[midway,below] {$dx$};
  \node[above=1pt] at ($(O)!0.55!(A)$) {$L_3$};
  \node[above left=0pt] at ($(A)!0.5!(W)$) {$L_4$};
\end{tikzpicture}
\caption{Esquema del mecanismo. El punto de salida del cable D está desplazado $dx$ respecto de la
vertical de la articulación O.}
\label{fig:mecanismo}
\end{figure}

El modelo parte de las siguientes hipótesis:
\begin{enumerate}[label=H\arabic*.]
  \item Régimen cuasiestático: a velocidad constante se desprecian las fuerzas de inercia.
  \item Cable inextensible y sin masa, y punto de salida D fijo.
  \item Se desprecia la masa de la barra frente a la masa suspendida.
  \item Las pérdidas por rozamiento no entran en el modelo estático: se tratan aparte en la
        sección~\ref{sec:mod-separacion}.
\end{enumerate}

\section{Cadena de transmisión}
\label{sec:mod-capstan}

Entre el servomotor y el tambor hay una reducción. La solución adoptada es una transmisión por
cable de tipo \emph{capstan drive}: el eje del motor lleva un cabrestante de radio pequeño sobre el
que se enrolla varias vueltas una cuerda sintética, cuyos extremos se anclan a una polea de radio
mayor solidaria al tambor. La relación de reducción es la de los radios,
\begin{equation}
  N = \frac{R_{\text{polea}}}{R_{\text{cabrestante}}},
  \label{eq:reduccion}
\end{equation}
con la salvedad de que el radio efectivo del cabrestante incluye el espesor de la cuerda y el paso
de la hélice con que se guía, de modo que conviene comprobarlo experimentalmente.

La cuerda no se anuda al cabrestante: se sujeta por rozamiento. La ecuación del cabrestante
relaciona las tensiones a uno y otro lado de un arrollamiento de ángulo total $\phi_w$ con el
coeficiente de rozamiento $\mu$,
\begin{equation}
  \frac{T_{\text{alta}}}{T_{\text{baja}}} \le e^{\mu\,\phi_w},
  \label{eq:capstan}
\end{equation}
un crecimiento exponencial con el número de vueltas que explica por qué con tres a cinco vueltas el
deslizamiento deja de ser un problema práctico.

Frente a un reductor de engranajes, esta transmisión no tiene holgura, tiene poca inercia reflejada
y es transparente al par, lo que importa aquí porque el par del accionamiento se usa como magnitud
de medida: cualquier holgura o rozamiento del reductor contaminaría precisamente la señal que se
quiere interpretar. A cambio, el recorrido es limitado —la cuerda sólo puede enrollarse y
desenrollarse— y exige retensado periódico. Para este banco la limitación no es tal: el mecanismo
trabaja entre dos finales de carrera y, como se verá, el recorrido completo son
\SI{0.209}{rev} de tambor, es decir, unos \SI{75}{\degree}.

La cuerda es de Dyneema (polietileno de ultra alto peso molecular). Se prefiere al cable de acero
porque admite radios de curvatura mucho menores sin fatigarse: un cable de acero exige una relación
entre el diámetro del tambor y el del cable elevada, y con radios pequeños falla por fatiga en pocas
horas. El punto débil del Dyneema es la fluencia —alargamiento permanente bajo carga mantenida—, que
se traduce en pérdida de tensión y error de posicionamiento, y que obliga a un tensor.
\completar{marca, diámetro y carga de rotura de la cuerda; radios del cabrestante y de la polea;
relación de reducción final medida}

\section{Longitud del cable}

Con $|OD| = \sqrt{L_2^2 + dx^2}$, la ley del coseno en el triángulo O--D--A da
\begin{equation}
  c(\alpha)^2 = |OD|^2 + L_3^2 - 2\,|OD|\,L_3\cos\alpha .
  \label{eq:cable}
\end{equation}
El coseno es monótono en $[0^\circ,180^\circ]$, de modo que la correspondencia entre $\alpha$ y $c$
es biunívoca: cada longitud de cable define un único ángulo, y al revés. Los extremos geométricos
son $c_{\min} = \big||OD|-L_3\big|$ con la barra plegada y $c_{\max} = |OD|+L_3$ con la barra
alineada. Despejando,
\begin{equation}
  \cos\alpha = \frac{|OD|^2 + L_3^2 - c^2}{2\,|OD|\,L_3}.
  \label{eq:angulo}
\end{equation}

\section{Tensión del cable}
\label{sec:mod-tension}

El momento del peso respecto a O es $mg(L_3+L_4)\cos\theta$. El brazo de la tensión es la distancia
de O a la recta D--A, que vale
\begin{equation}
  d = \frac{|OD|\,L_3\,\operatorname{sen}\alpha}{c}.
\end{equation}
Igualando momentos, $T\,d = mg(L_3+L_4)\cos\theta$, resulta
\begin{equation}
  \boxed{\;T = \frac{m\,g\,(L_3+L_4)\;c\,\cos\theta}{|OD|\;L_3\,\operatorname{sen}\alpha}\;}
  \qquad \theta = \beta - \alpha,\quad \beta = \arctan\frac{L_2}{-dx}.
  \label{eq:tension}
\end{equation}

Conviene detenerse en el caso particular $dx=0$, que fue el primer modelo del proyecto. Entonces
$\beta = 90^\circ$, $\theta = 90^\circ-\alpha$, $\cos\theta = \operatorname{sen}\alpha$, y los dos
factores se cancelan:
\begin{equation}
  T\big|_{dx=0} = \frac{m\,g\,(L_3+L_4)}{L_2\,L_3}\,c = K\,c ,
  \label{eq:tension0}
\end{equation}
es decir, \emph{la tensión sería proporcional a la longitud libre de cable}. Es un resultado
elegante, y por eso costó abandonarlo: con D desplazada la cancelación desaparece y la tensión deja
de ser proporcional al cable. Además, en $\alpha \to 180^\circ$ el brazo del cable se anula mientras
la pesa sigue haciendo momento, así que la tensión crece sin límite: ese punto es una singularidad
del mecanismo y no un punto de trabajo.

La tabla~\ref{tab:geometria} recoge los valores usados. No son cotas medidas, sino los parámetros
con los que trabaja la configuración del banco, ajustados sobre los propios ensayos
(sección~\ref{sec:res-ajuste}).

\begin{table}[H]
\centering
\caption{Parámetros del mecanismo empleados en el modelo \verificar{medir las cotas reales del
prototipo y rehacer el ajuste; de ellas dependen todas las cifras de este capítulo}.}
\label{tab:geometria}
\begin{tabular}{@{}lll@{}}
\toprule
Parámetro & Valor & Comentario \\
\midrule
$L_2$ & \SI{290}{\milli\metre} & altura de D sobre O \\
$dx$ & \SI{65}{\milli\metre} & desplazamiento horizontal de D \\
$|OD|$ & \SI{297.2}{\milli\metre} & línea O$\to$D, a \SI{12.6}{\degree} de la vertical \\
$L_3$ & \SI{95}{\milli\metre} & anclaje del cable en la barra \\
$L_4$ & \SI{110}{\milli\metre} & resto de barra hasta la masa \\
$m$ & \SI{2}{\kilo\gram} & masa suspendida \\
$\alpha_0$ & \SI{114}{\degree} & ángulo en B, donde arranca el ensayo \\
$c$ en B & \SI{346.9}{\milli\metre} & de~\eqref{eq:cable} \\
$c_{\min}$ -- $c_{\max}$ & \SIrange{202.2}{392.2}{\milli\metre} & topes geométricos \\
$r_0 \to r_1$ & \SIrange{45}{25}{\milli\metre} & radios de la caracola configurados \\
Barrido & \SI{0.3}{rev} de tambor & recorrido en que el radio pasa de $r_0$ a $r_1$ \\
$N$ & 5,5 & reducción motor:tambor \\
\bottomrule
\end{tabular}
\end{table}

\section{Par en el tambor}

Sea cual sea la forma del tambor, el cable sale por tangencia y su recta de acción queda a una
distancia $r(\varphi)$ del eje de giro. Esa distancia es a la vez el brazo del par y la velocidad
con que se paga cable por radián girado:
\begin{equation}
  \frac{\mathrm{d}s}{\mathrm{d}\varphi} = r(\varphi), \qquad
  s(\varphi) = \int_0^{\varphi} r(u)\,\mathrm{d}u, \qquad
  c(\varphi) = c_B - s(\varphi), \qquad
  \varphi = \frac{2\pi\,p}{N},
  \label{eq:cable-tambor}
\end{equation}
tomando como origen el punto B, donde arranca el ensayo, contando $\varphi$ positivo en el sentido
en que el ciclo avanza y siendo $p$ el avance del eje del motor en revoluciones. Que el brazo del
par y el radio de pago sean la misma función es lo que liga indisolublemente el perfil del tambor
con el recorrido: no se puede cambiar uno sin cambiar el otro. El par en el tambor y en el motor son
\begin{equation}
  \tau(\varphi) = T\big(c(\varphi)\big)\; r(\varphi), \qquad \tau_m = \frac{\tau}{N\,\eta},
  \label{eq:par}
\end{equation}
con $\eta$ el rendimiento de la transmisión.

\subsection{Perfiles de tambor}

Se manejan tres familias. El \textbf{cilindro} tiene $r(\varphi)=R$ constante; el par sigue
entonces exactamente a la tensión. La \textbf{caracola} es una espiral de Arquímedes, con radio
lineal en el giro entre $r_0$ y $r_1$ a lo largo de un barrido $\varphi_b$:
\begin{equation}
  r(\varphi) = r_0 + (r_1-r_0)\,\frac{\varphi}{\varphi_b}, \qquad 0 \le \varphi \le \varphi_b,
  \label{eq:espiral}
\end{equation}
y fuera de ese tramo el radio se queda en el extremo que corresponda.

El \textbf{tambor excéntrico} es un cilindro de radio $R$ cuyo eje de giro no pasa por su centro
geométrico, sino a una distancia $e$ de él (figura~\ref{fig:excentrico}). Es la forma que adopta el
diseño de este trabajo, y tiene una ventaja de fabricación decisiva: no hay ningún perfil que
mecanizar, basta taladrar el alojamiento del eje descentrado.

\begin{figure}[H]
\centering
\begin{tikzpicture}[scale=1.25,>=Stealth]
  \draw[violet,thick,fill=violet!7] (0,0) circle (1.6);
  \coordinate (C) at (0,0);
  \coordinate (Oe) at (-0.345,0.289);
  \draw[blue!70!black,thick] (1.6,1.5) -- (1.6,-1.9);
  \draw[->,thick,red!75!black] (1.6,-1.4) -- (1.6,-1.95) node[right] {$T$};
  \node[blue!70!black,right,font=\small] at (1.66,1.25) {cable};
  \fill (1.6,0.289) circle (1.1pt);
  \draw[dashed,gray] (C) -- (1.6,0);
  \node[gray,font=\small] at (0.8,-0.2) {$R$};
  \draw[->,green!45!black,thick] (Oe) -- (1.6,0.289);
  \node[green!45!black,font=\small,above] at (0.62,0.30) {$r(\varphi)$};
  \draw[->,thick] (C) -- (Oe);
  \node[font=\small,below right=-2pt] at (-0.17,0.145) {$e$};
  \fill (C) circle (1.3pt) node[below=2pt,font=\small] {C};
  \draw[line width=1pt] ($(Oe)+(-0.12,-0.12)$) -- ($(Oe)+(0.12,0.12)$);
  \draw[line width=1pt] ($(Oe)+(-0.12,0.12)$) -- ($(Oe)+(0.12,-0.12)$);
  \node[font=\small,left=3pt] at (Oe) {O$_t$};
  \draw[->,gray] (-1.05,0.95) arc (140:100:1.35);
  \node[gray,font=\small] at (-0.62,1.35) {$\varphi$};
\end{tikzpicture}
\caption{Tambor excéntrico: un cilindro de radio $R$ girando alrededor de un eje O$_t$ situado a una
distancia $e$ de su centro C. El brazo del cable es la distancia de O$_t$ a la recta de acción, y
varía de forma sinusoidal con el giro.}
\label{fig:excentrico}
\end{figure}

El cable sale tangente a la circunferencia de centro C y radio $R$, luego su recta de acción está a
distancia $R$ de C. Como C gira alrededor de O$_t$ describiendo un círculo de radio $e$, la
distancia de O$_t$ a esa recta es
\begin{equation}
  \boxed{\;r(\varphi) = R - e\cos(\varphi + \varphi_0)\;}
  \label{eq:excentrico}
\end{equation}
donde $\varphi_0$ es la fase de montaje, es decir, la orientación con que se cala el tambor sobre su
eje. Integrando~\eqref{eq:cable-tambor}, el cable pagado vale
\begin{equation}
  s(\varphi) = R\,\varphi - e\,\big[\operatorname{sen}(\varphi+\varphi_0) - \operatorname{sen}\varphi_0\big].
  \label{eq:excentrico-s}
\end{equation}

El radio recorre así una sinusoide de valor medio $R$ y amplitud $e$, entre $R-e$ y $R+e$, con
período una vuelta completa. De ahí se siguen dos propiedades que condicionan el diseño: el
recorrido de radios está acotado por la excentricidad, de modo que no se puede pedir cualquier
relación $r_{\max}/r_{\min}$; y el radio sólo es monótono durante media vuelta, así que el
mecanismo debe trabajar dentro de esa media vuelta. Ambas se cumplen holgadamente aquí, donde el
recorrido son \SI{75}{\degree}.

\subsection{Cota inferior del par de pico}

El trabajo que el tambor entrega a lo largo de un recorrido es el que gana la masa en energía
potencial, y no depende del perfil:
\begin{equation}
  W = \int_0^{\varphi_{\text{tot}}} \tau\,\mathrm{d}\varphi = \int_{c_{\text{fin}}}^{c_{\text{ini}}} T(c)\,\mathrm{d}c .
  \label{eq:trabajo}
\end{equation}
Como la integral del par está fijada, su máximo no puede ser menor que su media:
\begin{equation}
  \max_{\varphi}\tau(\varphi) \;\geq\; \frac{W}{\varphi_{\text{tot}}},
  \label{eq:cota}
\end{equation}
y la igualdad sólo se alcanza si el par es constante. \emph{El perfil de par constante es, por
tanto, el que minimiza el par de pico para un recorrido y un giro de tambor dados.} Este resultado
no depende de la geometría: vale igual con $dx=0$ que sin él.

\subsection{Perfil ideal de par constante}

Imponiendo $\tau = T(c)\,r = \tau_0$ y usando $\mathrm{d}c/\mathrm{d}\varphi = -r$:
\begin{equation}
  r(\varphi) = \frac{\tau_0}{T\big(c(\varphi)\big)}, \qquad
  \frac{\mathrm{d}c}{\mathrm{d}\varphi} = -\frac{\tau_0}{T(c)},
  \qquad \tau_0 = \frac{W}{\varphi_{\text{tot}}} .
  \label{eq:ideal}
\end{equation}
Con $dx=0$, donde $T=Kc$, la ecuación se integra en forma cerrada y da
$c(\varphi)=\sqrt{c_0^2-2(\tau_0/K)\varphi}$, con el radio inversamente proporcional al cable. Con
$dx\neq0$ no hay forma cerrada y la integración se hace numéricamente, que es como se ha calculado
la curva de la figura~\ref{fig:excentrico-comp}.

La regla práctica se mantiene en ambos casos: \emph{el radio debe ser grande donde la tensión es
pequeña y pequeño donde la tensión es grande}.

\subsection{¿Basta un tambor excéntrico?}
\label{sec:mod-basta}

El perfil ideal es una función cualquiera, que habría que mecanizar. El excéntrico sólo tiene dos
parámetros de forma, $e$ y $\varphi_0$, porque el tercero queda fijado: \emph{el tambor tiene que
pagar el mismo cable en el mismo giro}, ya que la carga recorre lo mismo. Imponiendo
$s(\varphi_{\text{tot}}) = S$ en~\eqref{eq:excentrico-s},
\begin{equation}
  R = \frac{S + e\,\big[\operatorname{sen}(\varphi_{\text{tot}}+\varphi_0) - \operatorname{sen}\varphi_0\big]}
           {\varphi_{\text{tot}}} .
  \label{eq:restriccion}
\end{equation}
Esta restricción no es un detalle de cálculo: sin ella el problema degenera, porque un tambor de
radio nulo no mueve la carga, entrega par constante cero y sale formalmente «óptimo».

La comparación se hace sobre el recorrido que los ensayos hacen realmente con carga en el cable
(sección~\ref{sec:res-campania}): un avance del motor de \num{0.65} a \SI{1.80}{rev}, que son
\SI{0.209}{rev} de tambor y llevan el cable de \num{316.4} a \SI{276.6}{\milli\metre}, es decir
$S = \SI{39.8}{\milli\metre}$. En ese tramo el modelo da una tensión que \emph{baja} de \num{44.2} a
\SI{35.5}{\newton}, porque el brazo del cable mejora más deprisa de lo que crece el momento del
peso, y un trabajo $W = \SI{1.58}{\joule}$, al que corresponde un par ideal
$\tau_0 = \SI{1.20}{\newton\metre}$.

Minimizando el rizado de par sobre $(e,\varphi_0)$ con $R$ dado por~\eqref{eq:restriccion} se
obtiene $R = \SI{32.8}{\milli\metre}$, $e = \SI{6.4}{\milli\metre}$ y
$\varphi_0 = \SI{27}{\degree}$, esto es, una excentricidad relativa de sólo $e/R = \num{0.19}$. El
resultado es el de la tabla~\ref{tab:perfiles}:

\begin{table}[H]
\centering
\caption{Perfiles de tambor sobre el recorrido con carga ($\varphi_{\text{tot}}=\SI{0.209}{rev}$ de
tambor, $S=\SI{39.8}{\milli\metre}$, $W=\SI{1.58}{\joule}$). Todos pagan el mismo cable en el mismo
giro.}
\label{tab:perfiles}
\begin{tabular}{@{}lcccc@{}}
\toprule
Perfil & $r$ inicial & $r$ final & $\tau$ mín. -- máx. & Rizado \\
       & (mm) & (mm) & (N·m) & $\frac{\max-\min}{\text{media}}$ \\
\midrule
Cilindro equivalente & 30,3 & 30,3 & 1,06 -- 1,35 & \SI{23.8}{\percent} \\
Caracola configurada & 41,2 & 27,9 & 0,88 -- 1,65 & \SI{64.5}{\percent} \\
\textbf{Excéntrico óptimo} & \textbf{27,1} & \textbf{34,2} & \textbf{1,197 -- 1,205} & \textbf{\SI{0.66}{\percent}} \\
Ideal (par constante) & 27,0 & 34,3 & 1,20 -- 1,20 & \SI{0}{\percent} \\
\bottomrule
\end{tabular}
\end{table}

La conclusión es la que da título a la sección: \emph{un excéntrico basta}. Su radio recorre de
\num{27.1} a \SI{34.2}{\milli\metre}, y el perfil ideal va de \num{27.0} a
\SI{34.3}{\milli\metre}; las dos curvas son indistinguibles a la escala del dibujo, y el rizado
residual es del \SI{0.66}{\percent} frente al \SI{0}{\percent} teórico. Dicho de otro modo: sobre un
recorrido de \SI{75}{\degree} el perfil ideal es, con excelente aproximación, un trozo de sinusoide,
y una sinusoide es exactamente lo que da un cilindro descentrado. Se obtiene el beneficio completo
de la compensación con una pieza que no requiere mecanizado de perfil.

Conviene subrayar también el contraste con la caracola \emph{tal y como está configurada}, que
empeora el reparto del par respecto a un simple cilindro (rizado del \SI{64.5}{\percent} frente al
\SI{23.8}{\percent}). La razón es que la tensión baja a lo largo del recorrido, de modo que para
compensarla el radio tendría que \emph{crecer} un \SI{26}{\percent}, y el perfil configurado lo hace
decrecer un \SI{32}{\percent}: está descrito —o montado— al revés respecto de lo que pide la
compensación.

\begin{figure}[H]
\centering
\includegraphics[width=\textwidth]{n_excentrico.pdf}
\caption{Izquierda: radio efectivo medido frente a las hipótesis de perfil
(sección~\ref{sec:res-contraste}). Derecha: par en el tambor para el cilindro equivalente, la
caracola configurada, el excéntrico óptimo y el perfil ideal; el excéntrico se superpone al ideal.}
\label{fig:excentrico-comp}
\end{figure}

\subsection{La pieza del prototipo}
\label{sec:mod-pieza}

El tambor que se ha llevado a fabricación no es el óptimo teórico de la
tabla~\ref{tab:perfiles}, sino una pieza con cotas redondeadas a lo que permitían el alojamiento del
rodamiento y el anclaje de la banda. Del modelo se toman dos radios sobre la misma diametral,
medidos desde el eje de giro: \SI{28.179}{\milli\metre} y \SI{56.00}{\milli\metre}
\verificar{confirmar que las dos cotas del CAD están sobre la misma recta que pasa por el eje; si no
lo están, $R$ y $e$ hay que obtenerlos del perfil completo}. De ahí,
\begin{equation}
  R = \frac{r_{\max}+r_{\min}}{2} = \SI{42.1}{\milli\metre}, \qquad
  e = \frac{r_{\max}-r_{\min}}{2} = \SI{13.9}{\milli\metre}, \qquad
  \frac{e}{R} = \num{0.33} .
\end{equation}

Al ser un tambor mayor que el óptimo, paga el cable en menos giro: el recorrido con carga pasa de
\SI{75.4}{\degree} a \SI{49.2}{\degree} de tambor (\SI{0.137}{rev}). Con el calado adecuado el
rizado de par resulta del \SI{2.8}{\percent}, con el radio recorriendo de \num{40.6} a
\SI{51.6}{\milli\metre} y el par entre \num{1.81} y \SI{1.86}{\newton\metre}. Sigue siendo un
resultado excelente ---casi un orden de magnitud mejor que el \SI{23.8}{\percent} de un cilindro---
aunque no alcance el \SI{0.66}{\percent} del óptimo. Manteniendo $R=\SI{42.1}{\milli\metre}$, la
excentricidad que minimizaría el rizado sería $e = \SI{10.0}{\milli\metre}$ ($e/R = \num{0.24}$),
que lo dejaría en el \SI{0.50}{\percent}: la pieza tiene algo más de excentricidad de la que este
mecanismo necesita.

\begin{table}[H]
\centering
\caption{Efecto del calado del tambor sobre su eje, con las cotas del prototipo.}
\label{tab:calado}
\begin{tabular}{@{}lcc@{}}
\toprule
Calado & $\varphi_0$ & Rizado de par \\
\midrule
Mejor & \SI{+84}{\degree} & \SI{2.8}{\percent} \\
Peor & --- & \SI{60.5}{\percent} \\
\bottomrule
\end{tabular}
\end{table}

Ese resultado merece destacarse porque es una instrucción de montaje, no un cálculo de diseño:
\emph{la misma pieza, calada en la posición equivocada, pasa de un rizado del \SI{2.8}{\percent} a
uno del \SI{60.5}{\percent}}, peor que no poner compensación ninguna. La razón es que el excéntrico
sólo compensa mientras su radio crece a lo largo del recorrido, y con media vuelta de diferencia en
el calado el radio decrece y suma su efecto al de la tensión en lugar de restarlo. El plano debe
acotar, por tanto, la posición angular del chavetero o de los taladros de arrastre respecto del
máximo de excentricidad, y el montaje debe verificarse antes de cargar la pesa.
\completar{incluir el plano de la pieza con la cota angular del calado y una vista del conjunto}

\subsection{Par exigido al motor}

La tabla~\ref{tab:parmotor} recoge el par de pico que cada solución exige al motor, repartiendo el
par del tambor por la reducción. El margen es amplio en todos los casos frente al par nominal del
servomotor, \SI{1.27}{\newton\metre}, lo que deja espacio para el rozamiento de la transmisión y
para el par de aceleración, que el modelo cuasiestático no recoge.

\begin{table}[H]
\centering
\caption{Par de pico exigido al motor, como fracción del par nominal (\SI{1.27}{\newton\metre}).}
\label{tab:parmotor}
\begin{tabular}{@{}lccc@{}}
\toprule
Tambor & $\tau$ de pico & \multicolumn{2}{c}{$\tau$ en el motor} \\
\cmidrule(lr){3-4}
 & en el tambor & $N=5{,}5$ & $N=\num{8.55}$ \\
\midrule
Excéntrico óptimo & \SI{1.205}{\newton\metre} & \SI{0.219}{\newton\metre} (\SI{17}{\percent}) & \SI{0.141}{\newton\metre} (\SI{11}{\percent}) \\
Excéntrico del prototipo & \SI{1.86}{\newton\metre} & \SI{0.338}{\newton\metre} (\SI{27}{\percent}) & \SI{0.218}{\newton\metre} (\SI{17}{\percent}) \\
Caracola configurada & \SI{1.650}{\newton\metre} & \SI{0.300}{\newton\metre} (\SI{24}{\percent}) & \SI{0.193}{\newton\metre} (\SI{15}{\percent}) \\
\bottomrule
\end{tabular}
\end{table}

El excéntrico del prototipo pide más par que el óptimo por una razón puramente geométrica: al ser un
tambor de mayor radio, el mismo trabajo se reparte en menos giro, de modo que el par sube en la
misma proporción en que baja el recorrido angular. Es un intercambio libre ---se compensa con la
reducción--- y no debe confundirse con el efecto de la compensación, que es lo que mide el rizado.

\section{Método experimental: separación del par y radio efectivo}
\label{sec:mod-separacion}

El accionamiento entrega el par en unidades internas de escala desconocida, y el par medido incluye
las pérdidas por rozamiento. Para comparar con el modelo sin depender de esa escala se recorre el
mismo punto en los dos sentidos, a la misma velocidad. Las pérdidas se oponen siempre al
movimiento, mientras que el par gravitatorio conserva el signo:
\begin{equation}
  |\tau_{\uparrow}(p)| = \tau_g(p) + \tau_f(p), \qquad
  |\tau_{\downarrow}(p)| = \tau_g(p) - \tau_f(p),
\end{equation}
\begin{equation}
  \tau_g = \frac{|\tau_{\uparrow}| + |\tau_{\downarrow}|}{2}, \qquad
  \tau_f = \frac{|\tau_{\uparrow}| - |\tau_{\downarrow}|}{2}.
  \label{eq:separacion}
\end{equation}
El mismo promedio aplicado a la señal de la célula elimina en primera aproximación el retardo de su
filtro, que desplaza en sentidos opuestos las curvas de ida y de vuelta.

Con $\tau_g$ ya limpio de pérdidas, la ecuación~\eqref{eq:par} permite estimar el \emph{radio
efectivo}
\begin{equation}
  r_{\text{ef}}(p) \;\propto\; \frac{\tau_g(p)}{T(p)},
  \label{eq:radio-efectivo}
\end{equation}
donde $T$ se toma del modelo, que no necesita la célula de carga: sólo la geometría y la masa. En un
tambor cilíndrico $r_{\text{ef}}$ sería constante; su forma es, por tanto, una medida directa del
perfil del tambor.

\subsection{Por qué no basta con medir la fuerza}
\label{sec:mod-ambiguedad}

Podría pensarse que basta con comparar la señal de la célula con la tensión del modelo. No es así, y
esto condiciona todo el capítulo de resultados. Como la célula no está calibrada, la comparación se
hace ajustando una recta $T \approx a\,(\text{cuentas}) + b$, y el coeficiente de determinación de
ese ajuste es el cuadrado del coeficiente de correlación:
\begin{equation}
  R^2 = \rho^2 .
\end{equation}
El valor de $R^2$ \emph{no cambia si se invierte el signo de $a$}. Una geometría que prediga una
tensión creciente y otra que la prediga decreciente pueden dar exactamente el mismo $R^2$ frente a
los mismos datos, y sólo difieren en el signo de la sensibilidad de la célula, que es precisamente
lo que no se conoce. La medida de fuerza, por sí sola, no puede decidir en qué sentido trabaja el
mecanismo.

\subsection{El mismo riesgo al comparar pares}
\label{sec:mod-metrica}

La advertencia anterior se extiende a la comparación de pares si se emplea una métrica ciega al
signo, y conviene explicitarlo porque es un error fácil de cometer. El par del accionamiento sí
tiene escala y signo conocidos, de modo que la comparación legítima ajusta \emph{una sola} constante
positiva $k$ —el paso de unidades internas a newton-metro— y mide la bondad como
\begin{equation}
  R^2 = 1 - \frac{\sum_i \big(\tau_{g,i} - k\,\tau_i^{\text{mod}}\big)^2}
                 {\sum_i \big(\tau_{g,i} - \bar\tau_g\big)^2}, \qquad k>0 .
  \label{eq:r2-par}
\end{equation}
Si en su lugar se usa $\rho^2$, un modelo que prediga exactamente la tendencia contraria a la medida
puntúa alto, porque $\rho^2$ es invariante frente al cambio de signo. La
sección~\ref{sec:res-contraste} muestra que la diferencia entre las dos métricas no es académica:
sobre estos datos una hipótesis pasa de $\rho^2=\num{0.81}$ —aparentemente buena— a
$R^2=\num{-6.10}$ —peor que predecir la media— según cuál se use.

\section{Limitaciones del modelo}

\begin{itemize}
  \item Los parámetros de la tabla~\ref{tab:geometria} no son cotas medidas, sino un ajuste sobre
        los ensayos; el propio capítulo~\ref{chap:resultados} muestra que no son consistentes con el
        par medido.
  \item El punto de salida del cable se desplaza con el radio de enrollamiento, lo que introduce un
        error de algunos milímetros en $L_2$ y $dx$.
  \item En el tambor excéntrico se supone que el cable trabaja en una sola capa y que su dirección
        de salida es fija. Con varias capas el radio efectivo crece con el espesor acumulado, y la
        dirección de salida gira, lo que altera~\eqref{eq:excentrico}.
  \item Se desprecian la elasticidad del cable y la masa de la barra. En una transmisión con cuerda
        sintética la fluencia del Dyneema hace que la primera hipótesis sea válida sólo a corto
        plazo, entre retensados.
  \item El rendimiento y las pérdidas se suponen iguales en los dos sentidos, en valor absoluto,
        para poder aplicar~\eqref{eq:separacion}.
  \item El modelo es estático: no describe el tramo inicial en que el cable está destensado, ni los
        transitorios de arranque e inversión.
\end{itemize}
